% ============================================================
% 实验报告通用模板 — lab-report-craft
% 文档类: ctexart (中文 article)
% 适用: 大学物理/近代物理/专业实验报告
% ============================================================
\documentclass[UTF8,a4paper,12pt]{ctexart}

% ---------- 页面与段落 ----------
\usepackage[top=2.45cm,bottom=2.45cm,left=2.55cm,right=2.55cm]{geometry}
\linespread{1.32}
\setlength{\parindent}{2em}
\tolerance=1200
\emergencystretch=3em

% ---------- 数学与符号 ----------
\usepackage{amsmath,amssymb,amsfonts}

% ---------- 图表 ----------
\usepackage{graphicx}
\usepackage{booktabs}
\usepackage{caption}
\captionsetup{font=small}
\usepackage{float}

% ---------- 列表缩进 ----------
\usepackage{enumitem}
\setlist[itemize]{leftmargin=2.4em}
\setlist[enumerate]{leftmargin=2.4em}

% ---------- 超链接 (打印友好) ----------
\usepackage{hyperref}
\hypersetup{colorlinks=true,linkcolor=black,citecolor=black,urlcolor=black}

% ---------- 章节样式 ----------
\ctexset{
  section={format=\large\bfseries, beforeskip=1.2ex, afterskip=0.8ex},
  subsection={format=\normalsize\bfseries, beforeskip=0.8ex, afterskip=0.4ex}
}

% ============================================================
% 封面信息 (请修改以下占位符)
% ============================================================
\newcommand{\reporttitle}{实验名称}
\newcommand{\studentname}{姓名}
\newcommand{\studentid}{学号}
\newcommand{\department}{院系}
\newcommand{\reportdate}{\today}

\begin{document}

% ---------- 封面 ----------
\begin{titlepage}
  \centering
  \vspace*{2cm}
  {\LARGE\bfseries 中国科学技术大学}\\[0.5cm]
  {\Large 实验报告}\\[2.5cm]
  {\Huge\bfseries \reporttitle}\\[3cm]
  \begin{tabular}{rl}
    \textbf{姓名:} & \studentname \\[0.3cm]
    \textbf{学号:} & \studentid \\[0.3cm]
    \textbf{院系:} & \department \\[0.3cm]
    \textbf{日期:} & \reportdate \\
  \end{tabular}
  \vfill
\end{titlepage}

% ---------- 摘要 ----------
\section*{摘要}
\addcontentsline{toc}{section}{摘要}

% 简明概括实验目的、方法、主要结果和结论 (200–300字)
% 例如：本实验利用...测量了...，得到...，与标准值...一致。

\vspace{0.5cm}
\noindent\textbf{关键词：}关键词1；关键词2；关键词3

% ---------- 1 实验目的 ----------
\section{实验目的}

\begin{itemize}
  \item 目的一
  \item 目的二
  \item 目的三
\end{itemize}

% ---------- 2 实验原理 ----------
\section{实验原理}

% 包含公式推导、物理图像描述
% 建议配合原理示意图 (Figure~\ref{fig:principle})

\begin{equation}
  % 主要公式
  \label{eq:main}
\end{equation}

% ---------- 3 实验装置与内容 ----------
\section{实验装置与内容}

\subsection{实验装置}

% 用表格列出主要仪器参数
\begin{table}[H]
  \centering
  \caption{主要实验仪器}
  \label{tab:instruments}
  \begin{tabular}{lll}
    \toprule
    仪器名称 & 型号 & 主要参数 \\
    \midrule
    % 填写仪器信息
    \bottomrule
  \end{tabular}
\end{table}

\subsection{实验步骤}

\begin{enumerate}
  \item 步骤一
  \item 步骤二
  \item 步骤三
\end{enumerate}

% ---------- 4 实验结果与分析 ----------
\section{实验结果与分析}

\subsection{结果一 (如：数据处理)}

% 插入图表，caption 说明数据来源和物理意义
\begin{figure}[H]
  \centering
  \includegraphics[width=0.85\textwidth]{figures/figure1.png}
  \caption{图表说明。数据来源：自主测量。}
  \label{fig:result1}
\end{figure}

\subsection{误差分析}

% 定量分析误差来源：统计误差、系统误差、模型误差等

% ---------- 5 结论 ----------
\section{结论}

% 总结主要定量结果，给出不确定度
% 若有多个数据集，明确区分自主测量与参考数据
% 指出实验的局限性与改进方向

% ---------- 附录 ----------
\newpage
\appendix
\section{思考题解答}

% 按题目顺序给出解答

\clearpage
\section{原始数据}

% 手写记录照片或完整数据表格

% ---------- 参考文献 ----------
\clearpage
\begin{thebibliography}{99}
\bibitem{ref1} 作者, ``标题,'' \textit{期刊名} \textbf{卷}, 页码 (年份).
\bibitem{ref2} 机构名, ``报告标题,'' URL.
\end{thebibliography}

\end{document}
